\chapter{Hỏi Cho Giờ Chủ Nhiệm}
\chapterintro{Muốn lớp nói thật, câu hỏi phải vừa mở vừa an toàn}{Giờ chủ nhiệm rất cần những câu hỏi giúp lớp soi mình mà không bật chế độ phòng thủ.}

\section{Tuần này lớp mình nên tự khen điều gì trước?}
\begin{cauhoibox}{Câu hỏi}
Tuần này lớp mình nên tự khen điều gì trước?
\end{cauhoibox}
\begin{taisaocoich}
Mở bằng điểm tốt giúp lớp dễ nói thật hơn ở phần cần sửa sau đó.
\end{taisaocoich}
\begin{goiybien}
Dùng cực hợp đầu giờ sinh hoạt để tránh không khí kiểm điểm nặng ngay từ đầu.
\end{goiybien}
\ngancach

\section{Nếu lớp mình chỉ sửa được một điều trong tuần tới, nên sửa điều gì?}
\begin{cauhoibox}{Câu hỏi}
Nếu lớp mình chỉ sửa được một điều trong tuần tới, nên sửa điều gì?
\end{cauhoibox}
\begin{taisaocoich}
Câu hỏi giới hạn vào một điều khiến lớp thực tế hơn, tránh nói quá nhiều thứ rồi không thay đổi được gì.
\end{taisaocoich}
\begin{goiybien}
Sau khi chọn, viết điều đó lên bảng và nhắc lại ở cuối buổi sinh hoạt.
\end{goiybien}
\ngancach

\section{Lớp mình cần thêm điều gì để thấy vui khi đi học hơn?}
\begin{cauhoibox}{Câu hỏi}
Lớp mình cần thêm điều gì để thấy vui khi đi học hơn?
\end{cauhoibox}
\begin{taisaocoich}
Câu hỏi này đưa giờ chủ nhiệm về đúng tinh thần xây lớp, không chỉ sửa lỗi. Nó thường kéo được nhiều câu trả lời thật và dễ thương.
\end{taisaocoich}
\begin{goiybien}
Hợp để làm cuối tháng hoặc sau khi lớp vừa trải qua một giai đoạn mệt.
\end{goiybien}
\ngancach

\section{Nếu lớp mình có một điều muốn nói thật hôm nay, đó là gì?}
\begin{cauhoibox}{Câu hỏi}
Nếu lớp mình có một điều muốn nói thật hôm nay, đó là gì?
\end{cauhoibox}
\begin{taisaocoich}
Giờ chủ nhiệm cần những câu hỏi đủ an toàn để lớp mở lòng. Câu này mở rộng nhưng vẫn gọn, giúp kéo ra điều thật mà lớp đang giữ trong lòng.
\end{taisaocoich}
\begin{goiybien}
Nếu lớp ngại nói công khai, có thể cho viết giấy rồi thầy cô đọc lại theo nhóm ý.
\end{goiybien}
\ngancach

\section{Tuần tới em muốn bước vào lớp với cảm giác gì hơn?}
\begin{cauhoibox}{Câu hỏi}
Tuần tới em muốn bước vào lớp với cảm giác gì hơn?
\end{cauhoibox}
\begin{taisaocoich}
Câu hỏi này kéo giờ chủ nhiệm về phía xây lớp, không chỉ sửa lỗi. Nó thường tạo ra những câu trả lời rất thật mà thầy cô có thể dùng để chỉnh cách tổ chức lớp.
\end{taisaocoich}
\begin{goiybien}
Hợp để chốt cuối buổi sinh hoạt bằng một hướng đi tích cực thay vì chỉ dừng ở bản tổng kết.
\end{goiybien}
